% Stand 24. Juli 2026
\subsubsection{Kriging-Interpolation mit anisotropen Kovarianz-Modellen}
\label{KapitelDatenauswertung_Ergebnisse_Kriging}
Wie im vorherigen Kapitel \ref{KapitelDatenauswertung_Ergebnisse_Variogramme} 
ausgeführt, lässt sich auf der Grundlage der vorhandenen Messungen
keine belastbare Aussage hinsichtlich einer möglichen Anisotropie der 
Feinstaubdaten treffen. \\
Die in diesem Kapitel vorgestellte Methode versteht sich als ein explorativer
Ansatz, um \textit{Indizien} für oder gegen die Anisotropie-Vermutung herauszuarbeiten.
Zu diesem Zweck wird das folgende Vorgehen durchgeführt: 
\begin{itemize}[topsep=2pt,parsep=1pt,itemsep=0.5pt]
\item[1) ] bei schrittweiser Variation der Anisotropie mit $\Delta \lambda = 1$, wobei $1 \leq \lambda \leq 4$, sowie
\item[2) ] für drei verschiedene Wert des Sill $\mathrm{S}$, 
d.\,h.\,$\mathrm{S} \in \lbrace \mathrm{S}_{\mathrm{min}},\,\mathrm{S}_{\mathrm{mid}},\,\mathrm{S}_{\mathrm{max}}\rbrace$ und
\item[3) ] mit einer Winkel-Schrittweite von 
$\Delta\,\vartheta= 20^\circ$, wobei 
$0^\circ \leq \vartheta < 360^\circ$
\item[4) ] wird ein anisotropes Kovarianz-Modell definiert 
und damit für die beobachteten Werte die 
Kriging-Gleichungen gelöst. 
Insbesondere findet somit \textit{keine} Anpassung der
Modell-Parameter an 
die richtungsabhängig geschätzte Semivarianz statt, da sich
hierbei -- wie in Kapitel
\ref{KapitelDatenauswertung_Ergebnisse_Variogramme} ausgeführt -- ja keine verwertbaren Ergebnisse zeigen.
\item[5) ] Um die Güte der anisotropen Modell-Wahl zu 
evaluieren, wird
ein Leave-One-Out-Cross-Validation (\textit{kurz:} LOO-CV) 
Ansatz  
(vgl.\,z.\,B.\,\cite[Kap.\,3.8.4]{Murphy2023}, 
\cite[Kap.\,5.1.2]{James2023}) gewählt: 
\begin{itemize}[topsep=2pt,parsep=1pt,itemsep=0.5pt]
\item[a) ] Aus den $n$ Beobachtungen werden 
$(n-1)$ ausgewählt und mit deren Hilfe die 
Kriging-Gewichte bestimmt.
\item[b) ] Für die eine, noch nicht verwendete 
Beobachtung wird eine Vorhersage erzeugt und für diese die
quadrierte Differenz $e_j$ zur Beobachtung berechnet.
\item[c) ] Dieser Ablauf wird für alle 
$n$ möglichen Aufteilungen in Trainings- und Test-Menge
durchgeführt. 
\item[d) ] Damit berechnet sich der Fehler $e_\mathrm{MSE}$ der Cross-Validation:
\begin{equation*}
e_\mathrm{MSE} = \frac{1}{n} \sum_{j=1}^n e_j.
\end{equation*}
%\item[e)] Fehler $\sigma_\mathrm{MSE}$ der Cross-Validation:
%\begin{equation*}
%\sigma^2_\mathrm{MSE} = \frac{1}{n} \sum_{j=1}^n 
%\left( e_j - e_{\mathrm{MSE}} \right)^2.
%\end{equation*}
\end{itemize}
\end{itemize}

Da die Parameter des anisotropen Kovarianz-Modells \textit{nicht} 
mithilfe der empirischen Semivarianz anhand der Messdaten geschätzt 
werden können, ist das skizzierte Vorgehen methodisch eher als eine 
Art plausibilisiertes Raten zu bezeichnen und stellt 
sicher \textit{keine} saubere Analyse dar.\\

Es erweist sich, dass die einzelnen Werte der 
$\lbrace e_j\rbrace_{j \in [n]}$ typischerweise eine
Spannbreite von mehreren Größenordnungen abdecken, 
wobei die größten Werte genau dann auftreten, wenn
für ungewöhnlich große Beobachtungen der 
Feinstaub-Konzentration eine Vorhersage gemäß Schritt b) 
bestimmt wird und die Menge der korrespondierenden 
Trainingsdaten nach a) überwiegend sehr viel kleinere Feinstaub-Werte
beinhaltet (vgl.\,Abb.\,(\ref{Abb_Datenauswertung_P25_LOO-CVNov2025_Ansicht_e_j}) 
im Anhang). Diese nur wenigen, aber sehr großen Werte $e_j$ ziehen 
den als Mittelwert kalkulierten Fehler $e_{\mathrm{MSE}}$
nach oben und führen zu großen, durch die Standardabweichung 
quantifizierten Schwankungen. Des Weiteren musste mit $c_0 = 0.2$ dem Modell 
ein Nugget-Effekt hinzugefügt werden, da andernfalls die Kriging-Interpolation 
teilweise zu negativen Werten führt.\\

Die Ergebnisse dieser Analyse für den 
November 2025 sind für ein
festes Anisotropie-Verhältnis bzw.\,einen festen Sill in 
den 
Abb.\,(\ref{Abb_Datenauswertung_P25_LOO-CVNov2025_Lambda30}) 
und 
(\ref{Abb_Datenauswertung_P25_LOO-CVNov2025_Sill40}) 
wiedergegeben.

\begin{figure}[H] 
\centering
\includegraphics[scale = 0.45]{Datenauswertung_P25_LOO-CVNov2025_Lambda30} 
\caption{\small{Ergebnis der Leave-One-Out Cross-Validation
für ein fixes Anisotropie-Verhältnis von $\lambda=3$
basierend auf den gemittelten 
$\mathrm{P}_{2.5}$-Messwerten für November 2025}}
\label{Abb_Datenauswertung_P25_LOO-CVNov2025_Lambda30}
\end{figure} 

\begin{figure}[H] 
\centering
\includegraphics[scale = 0.45]{Datenauswertung_P25_LOO-CVNov2025_Sill40} 
\caption{\small{Ergebnis der Leave-One-Out Cross-Validation
für einen Sill von $\mathrm{S}=4$
basierend auf den gemittelten 
$\mathrm{P}_{2.5}$-Messwerten für November 2025}}
\label{Abb_Datenauswertung_P25_LOO-CVNov2025_Sill40}
\end{figure} 

Da nur das Anisotropie-Verhältnis $\lambda$ variiert wird,
muss für eine der beiden Halbachsen der Ellipse ein Wert
vorgegeben werden. Hierfür wurde die kleine Halbachse
für alle drei Datensätze auf den maximalen Wert der 
paarweisen Distanz $\mathrm{R}_{max}$ der 
Sensoren-Positionen gesetzt. Folglich liegen
die Werte der großen Halbachse bei etwa $30-60\,km$.\\

Aus beiden Abbildungen ersieht man:
\begin{itemize}[topsep=2pt,parsep=1pt,itemsep=0.5pt]
\item[$\circ$ ] den geringsten Prognosefehler
$e_{\mathrm{MSE}}$ weisen die konstruierten anisotropen 
Kovarianz-Modelle auf, die eine Drehung der Koordinaten in
Richtung $60^\circ$ (oder gleichbedeutend 
$240^\circ$) beinhalten. 
Diese Beobachtung ist unabhängig vom
gewählten Wert für Anisotropie und Sill.
\item[$\circ$ ] Tendenziell verringert sich mit steigender 
Anisotropie der Fehler $e_{\mathrm{MSE}}$.
\item[$\circ$ ] Tendenziell verringert sich mit abnehmendem Sill 
der Fehler $e_{\mathrm{MSE}}$.
\end{itemize}

\vspace*{4pt}
Auch die analogen Auswertungen der Messungen aus dem Juli 2025 
in 
Abb.\,(\ref{Abb_Datenauswertung_P25_LOO-CVJul2025_Sill08}) 
und dem April 2026 in 
Abb.\,(\ref{Abb_Datenauswertung_P25_LOO-CVApr2026_Sill10})
weisen den geringsten Prognose-Fehler für das Kovarianz-Modell auf, 
welches eine Drehung um etwa $40^\circ$ bzw.\,$180^\circ$ impliziert.
Hierbei sind diese Vorzugsrichtungen wiederum unabhängig vom
gewählten Wert für Anisotropie und Sill. \\

Für alle drei der monatlich gemittelten Datensätze, dargestellt in den 
Abbildungen 
(\ref{Abb_Datenauswertung_P25_LOO-CVNov2025_Sill40}),
(\ref{Abb_Datenauswertung_P25_LOO-CVJul2025_Sill08}) und
(\ref{Abb_Datenauswertung_P25_LOO-CVApr2026_Sill10}),
 liegt der Fehler
$e_{\mathrm{MSE}}$ des isotropen ($\lambda =1$) Modells 
(\textcolor{Green}{\rule[1mm]{3mm}{0.5mm}}) oberhalb des Fehlers 
für das anisotrope Modell mit der fehlerminimalen Drehung. Das 
isotrope Modell beschreibt die räumliche Verteilung der Feinstaubmessungen 
somit \textit{weniger gut}.\\

\begin{figure}[H] 
\centering
\includegraphics[scale = 0.45]{Datenauswertung_P25_LOO-CVJul2025_Sill08} 
\caption{\small{Ergebnis der Leave-One-Out Cross-Validation
für einen Sill von $\mathrm{S}=0.8$
basierend auf den gemittelten 
$\mathrm{P}_{2.5}$-Messwerten für Juli 2025}}
\label{Abb_Datenauswertung_P25_LOO-CVJul2025_Sill08}
\end{figure} 

\begin{figure}[H] 
\centering
\includegraphics[scale = 0.45]{Datenauswertung_P25_LOO-CVApr2026_Sill10} 
\caption{\small{Ergebnis der Leave-One-Out Cross-Validation
für einen Sill von $\mathrm{S}=1.0$
basierend auf den gemittelten 
$\mathrm{P}_{2.5}$-Messwerten für April 2026}}
\label{Abb_Datenauswertung_P25_LOO-CVApr2026_Sill10}
\end{figure} 

%Für alle drei der monatlich gemittelten Datensätze, dargestellt in den 
%Abbildungen 
%(\ref{Abb_Datenauswertung_P25_LOO-CVNov2025_Sill40}),
%(\ref{Abb_Datenauswertung_P25_LOO-CVJul2025_Sill08}) und
%(\ref{Abb_Datenauswertung_P25_LOO-CVApr2026_Sill10})
% liegt der Fehler
%$e_{\mathrm{MSE}}$ des isotropen ($\lambda =1$) Modells 
%(\textcolor{Green}{\rule[1mm]{3mm}{0.5mm}}) oberhalb des Fehlers 
%für das anisotrope Modell mit der fehlerminimalen Drehung. Das 
%isotrope Modell beschreibt die räumliche Verteilung der Feinstaubmessungen 
%somit \textit{weniger gut}.\\
Die Frage nach der statistischen Signifikanz ist nicht einfach zu klären: der 
naheliegend erscheinende t-Test für gepaarte Stichproben
 (vgl.\,z.\,B.\,\cite[S.\,263]{Henze2019}) 
würde die
paarweisen Differenzen der Fehler von isotroper und
 anisotroper Modellierung betrachten, diese als stochastisch unabhängig annehmen
und davon ausgehend einen gewöhnlichen t-Test ableiten. 
Da bei der Berechnung des LOO-CV aber dieselben räumlich 
verteilten Messgrößen eingehen, dürfte diese Annahme kaum erfüllt sein. 
Der auch im Hinblick auf andere Fragestellungen  
(vgl.\,Kapitel \ref{KapitelDatenauswertung_Diskussion}) 
geeignete Permuationstest könnte hier Klarheit bringen.\\
   
Der Vollständigkeit halber sind in 
Abb.\,(\ref{Abb_Datenauswertung_P25_KrigingMean_Winkel40_Sill4x0_Lambda3x0_Nov2025})
und
(\ref{Abb_Datenauswertung_P25_KrigingVariance_Winkel40_Sill4x0_Lambda3x0_Nov2025})
farbkodiert der Erwartungswert und die
Varianz der Kriging-Interpolation für ein konstruiertes,
anisotropes Kovarianz-Modell mit 
($\vartheta=40^\circ,\,\lambda = 3,\,\mathrm{S}=4)$ dargestellt\footnote{Der
Kriging-Algorithmus ist bei nicht-verschwindendem Nugget 
\textit{kein} Interpolator mehr. Dies erkennt man in der Abbildung 
prominent anhand der rot eingezeichneten maximalen Beobachtung.}. 

\begin{figure}[H] 
\centering
\includegraphics[scale = 0.45]{Datenauswertung_P25_KrigingMean_Winkel40_Sill4x0_Lambda3x0_Nov2025} 
\caption{\small{Erwartungswert der Kriging-Interpolation
für $\vartheta=40^\circ$, $\lambda=3$ und $\mathrm{S}=4$
basierend auf den gemittelten 
$\mathrm{P}_{2.5}$- Messwerten für November 2025}}
\label{Abb_Datenauswertung_P25_KrigingMean_Winkel40_Sill4x0_Lambda3x0_Nov2025}
\end{figure} 

\begin{figure}[H] 
\centering
\includegraphics[scale = 0.45]{Datenauswertung_P25_KrigingVariance_Winkel40_Sill4x0_Lambda3x0_Nov2025} 
\caption{\small{Varianz der Kriging-Interpolation
für $\vartheta=40^\circ$, $\lambda=3$ und $\mathrm{S}=4$
basierend auf den gemittelten 
$\mathrm{P}_{2.5}$- Messwerten für November 2025}}
\label{Abb_Datenauswertung_P25_KrigingVariance_Winkel40_Sill4x0_Lambda3x0_Nov2025}
\end{figure} 

%Die Anisotropie des zugrundeliegenden Kovarianz-Modells erkennt 
%man in 
%Abb.\,(\ref{Abb_Datenauswertung_P25_KrigingMean_Winkel40_Sill0x8_Lambda3x5_Nov2025_Lambda_fix}) 
%recht deutlich in der oberen linken Ecke anhand der ellipsoiden, 
%rot gefärbten Flächen.


